\cleardoublepage
\addtocontents{toc}{\vspace{2em}} % Add a gap in the Contents, for aesthetics
\appendix
\chapter{Dataset file format}
\label{app:dataset-format}

The dataset described in \Cref{sec:origin-structure} was stored as a single \texttt{.parquet} file. \Cref{tab:dataset-columns} details the exact column layout of that file, as used by the Python code, available at \href{https://github.com/claudiocaaruso/VibraAI}{\textcolor{blue}{\underline{this GitHub repository}}}.

\begin{table}[H]
    \centering
    \begin{tabularx}{\textwidth}{|l|l|X|}
        \hline
        \textbf{Column} & \textbf{Format / Range} & \textbf{Description} \\
        \hline
        \texttt{Sample\_ID} & e.g. \texttt{1234\_567} & The first 4 digits identify the patient, the
                              last 3 the sample from that patient. \\
        \hline
        \texttt{Map\_ID}    & e.g. \texttt{Map\_1} & Acquisition map; 2 to 5 maps were
                              analysed per sample. \\
        \hline
        \texttt{x} & $x\in [0, 119]$ & Pixel coordinates on the horizontal axis\\
         \hline
        \texttt{y} & $y\in [0, 119]$ & Pixel coordinates on the vertical axis\\
        \hline
        \texttt{band\_}$n$  & $n \in [0, 482]$ & Raman intensity at the $n$-th spectral band,
                              for the 483 Raman shifts ($\mathrm{cm}^{-1}$). \\
        \hline
        \texttt{Label}      & integer & Tissue-class identifier of the pixel (see \Cref{sec:labels}). \\
        \hline
    \end{tabularx}
    \caption{Column layout of the dataset's \texttt{.parquet} file.}
    \label{tab:dataset-columns}
\end{table}
